%%%%%%%%%%%%%%%%%%%%%%%%%%%%%%%%%%%%%%%%%%%%%%%%%%%%%%%%%%%%%%%%%%%%%%
% Lớp 11 | Xác suất | Biến cố giao và quy tắc nhân xác suất
% Dạng 1D9?1-1: Câu hỏi lý thuyết
% Số câu: 11  (sinh tự động bởi tach_cau.py; từ nay thầy sửa trực tiếp file này)
%%%%%%%%%%%%%%%%%%%%%%%%%%%%%%%%%%%%%%%%%%%%%%%%%%%%%%%%%%%%%%%%%%%%%%

%%%%%%%%%%%%%%%% TRẮC NGHIỆM 4 PHƯƠNG ÁN %%%%%%%%%%%%%%%%

% nguồn: Ok_TracNghiemLop11_Lan1
\begin{ex}%[1D9N1-1]%[ma:ef718e90]
	Cho hai biến cố $A$,  $B$ độc lập. Mệnh đề nào sau đây \textbf{sai}?
	 \choice
		{$\overline{A}$, $B$ là độc lập} 
		 {$A$, $\overline{B}$ là độc lập}
		{$\overline{A}$, $\overline{B}$  là độc lập}
		 {\True $A$, $\overline{A}$ là độc lập}
	 \loigiai{
		Mệnh đề sai là \lq\lq$A$,  $\overline{A}$ là độc lập\rq\rq.
	} 
\end{ex}

% nguồn: Ok_TracNghiemLop11_Lan1
\begin{ex}%[1D9N1-1]%[ma:730ee932]
	Cho $A$, $\overline{A}$ là hai biến cố đối nhau trong cùng một phép thử $T$; xác suất xảy  ra biến cố $A$ là $\dfrac{1}{3}$ . Xác suất để xảy ra biến cố $\overline{A}$ là
	\choice 
		 {$\mathrm{P}\left(\overline{A}\right)=1$}
		{\True  $\mathrm{P}\left(\overline{A}\right)=\dfrac{2}{3}$}
		 {$\mathrm{P}\left(\overline{A}\right)=\dfrac{1}{2}$}
		 {$\mathrm{P}\left(\overline{A}\right)=\dfrac{1}{3}$}
	\loigiai{ 
		Vì $A$, $\overline{A}$ là hai biến cố đối  nhau trong cùng một phép thử $T$ nên ta có
		\[\mathrm{P}(A)+\mathrm{P}\left(\overline{A}\right)=1 \Rightarrow  \mathrm{P}\left(\overline{A}\right)=1-\mathrm{P}(A)=1-\dfrac{1}{3}=\dfrac{2}{3}.\]
	 }
\end{ex}

% nguồn: Ok_TracNghiemLop11_Lan1
\begin{ex}%[1D9N1-1]%[ma:86f2580f]
	Cho $A$, $B$ là hai  biến cố liên quan đến một phép thử có hữu hạn các kết quả đồng khả năng xuất hiện. Xét $A$ và $B$ là hai biến cố độc lập. Mệnh đề nào sau đây \textbf{sai}?
	 \choice
		 {$A$ và $\overline{B}$ là hai biến cố độc lập}
		 {\True $\mathrm{P}\left(\overline{A}B\right)=\mathrm{P}(A)\cdot \mathrm{P}(B)$}
		{$\mathrm{P}(AB)=\mathrm{P}(A)\cdot  \mathrm{P}(B)$}
		 {$\mathrm{P}\left(\overline{A}B\right)=\left(1-\mathrm{P}(A)\right)\cdot \mathrm{P}(B)$}
	\loigiai{ 
		$\mathrm{P}\left(\overline{A}B\right)=\mathrm{P}(A)\cdot \mathrm{P}(B)$ là mệnh đề sai vì $\mathrm{P}\left(\overline{A}B\right)=\mathrm{P}(\overline{A})\cdot  \mathrm{P}(B)$.
	} 
\end{ex}

% nguồn: Ok_TracNghiemLop11_Lan1
\begin{ex}%[1D9N1-1]%[ma:f9c53f67]
	Một hộp có 52 chiếc thẻ cùng loại, mỗi thẻ được ghi một trong các số tự nhiên, hai thẻ khác nhau thì ghi hai số khác nhau. Rút ngẫu nhiên 1 chiếc thẻ trong hộp. Xét biến cố A: \lq\lq Số xuất hiện trên thẻ được rút ra là số lẻ\rq\rq\,và biến cố B: \lq\lq Số xuất hiện trên  thẻ được rút ra là số nguyên tố\rq\rq.\\
	 Chọn câu \textbf{sai} khi phát biểu về biến cố AB.
	\choice 
		{\True \lq\lq Số xuất hiện trên thẻ là  số lẻ hoặc số nguyên tố\rq\rq}
		{\lq\lq Số xuất hiện trên  thẻ là số nguyên tố lẻ\rq\rq }
		{\lq\lq  Số xuất hiện trên thẻ là số nguyên tố khác $2$\rq\rq }
		{\lq\lq Số xuất  hiện trên thẻ vừa là số lẻ vừa là số nguyên tố\rq\rq}
	\loigiai{ 
	Dựa vào định nghĩa có \lq\lq Số xuất hiện trên thẻ là số lẻ hoặc  số nguyên tố\rq\rq\, là phát biểu sai về AB.}
\end{ex}

% nguồn: Ok_TracNghiemLop11_Lan1
\begin{ex}%[1D9N1-1]%[ma:bb3c96db]
	Tung  một đồng xu cân đối và đồng chất hai lần liên tiếp. Xét các biến cố:\\
	X:  \lq\lq Đồng xu xuất hiện mặt sấp ở lần gieo thứ nhất\rq\rq .\\
	Y: \lq\lq Đồng  xu xuất hiện mặt ngửa ở lần gieo thứ hai\rq\rq .\\
	Hai biến cố X  và Y là hai
	 \choice
		 {biến cố đối}
		{\True biến cố độc lập} 
		 {biến cố giao}
		{biến cố  hợp}
	 \loigiai{
	Dựa vào  định nghĩa.}
\end{ex}

% nguồn: Ok_TracNghiemLop11_Lan1
\begin{ex}%[1D9N1-1]%[ma:2aff9bad]
	Cặp biến cố A và B được gọi là  độc lập nếu
	 \choice
		{Việc biến cố này xảy ra hay không  ảnh hưởng tới xác suất xảy ra của biển cố kia}
		{\True Việc xảy ra hay không xảy ra của biến cố này không ảnh hưởng tới xác suất xảy ra của  biến cố kia}
		{A hoặc  B xảy ra}
		{A và B cùng  xảy ra}
	\loigiai{ 
	Dựa  vào định nghĩa.}
\end{ex}

% nguồn: Ok_TracNghiemLop11_Lan1
\begin{ex}%[1D9N1-1]%[ma:c1eff063]
	Một hộp chứa $21$ tấm thẻ cùng loại được đánh số từ $1$ đến $21$. Chọn ra ngẫu nhiên $1$ thẻ  từ hộp. Gọi $A$ là biến cố \lq\lq Số ghi trên thẻ được chọn chia hết cho $2$\rq\rq, $B$ là biến cố \lq\lq Số ghi trên thẻ được chọn chia hết cho $3$\rq\rq.
	Chọn khẳng định \textbf{đúng} trong các khẳng định  sau:
	 \choice
		{Hai biến cố $A$  và $B$ độc lập}
		{Hai biến cố $A$ và $B$  xung khắc}
		{Hai biến  cố $A$ và $\overline{B}$ đối nhau}
		{\True  Hai biến cố $A$ và $B$ không độc lập}
	 \loigiai{
		 Ta có $A=\{2;4;6;8;10;12;14;16;18;20\}$, $B=\{3;6;9;12;15;18;21\}$.\\
		Suy ra  $P(A)=\dfrac{10}{21}$; $P(B)=\dfrac{7}{21}=\dfrac{1}{3}$ và $P(AB)=\dfrac{3}{21}=\dfrac{1}{7}$.\\
	Ta có $P(AB)=\dfrac{1}{7}\neq \dfrac{10}{21}\cdot \dfrac{1}{3}=\dfrac{10}{63}=P(A)P(B)$ nên hai biến cố $A$ và $B$ không  độc lập.}
\end{ex}

% nguồn: Ok_TracNghiemLop11_Lan1
\begin{ex}%[1D9N1-1]%[ma:a1747e06]
	Cho $A$, $B$ là hai biến cố liên quan đến một phép thử có hữu hạn các kết quả  đồng khả năng xuất hiện. Nếu $A\cap B=\varnothing$ thì $A$ và $B$ là hai biến cố
	 \choice
		{độc  lập}
		 {giao}
		{\True xung khắc} 
		 {đối nhau}
	 \loigiai{
		 Nếu $A\cap B=\varnothing$ thì $A$ và $B$ là hai biến cố xung khắc.
	} 
\end{ex}

% nguồn: Ok_TracNghiemLop11_Lan1
\begin{ex}%[1D9N1-1]%[ma:37e10314]
	Chọn ngẫu nhiên một học sinh trong một trường. Xét 2  biến cố sau:\\
	P: \lq\lq Học sinh đó bị cận thị\rq\rq  .\\
	Q:  \lq\lq Học sinh đó học giỏi Toán\rq\rq .\\
	Nêu nội dung  của các biến cố P giao Q là
	 \choice
		{Biến cố PQ xảy ra khi học sinh đó chỉ bị cận thị hoặc chỉ  học giỏi môn Toán}
		{Biến cố PQ xảy ra khi học sinh đó bị cận thị  hoặc học giỏi môn Toán}
		{\True Biến cố PQ xảy ra khi học sinh đó vừa bị cận thị vừa học giỏi  môn Toán}
		{Biến cố PQ xảy ra khi học sinh đó không bị cận thị và không  học giỏi môn Toán}
	 \loigiai{
	Dựa vào  định nghĩa.}
\end{ex}

% nguồn: Ok_TracNghiemLop11_Lan1
\begin{ex}%[1D9N1-1]%[ma:462b68ec]
	Một hộp đựng 5 viên bi màu đỏ và 6 viên bi màu xanh có cùng kích thước và khối lượng. Bạn Hưng lấy ngẫu nhiên một viên bi và không trả lại vào hộp. Tiếp theo bạn An lấy ngẫu nhiên một  viên bi từ hộp đó. Xét biến cố A: \lq\lq Hưng lấy được viên bi màu đỏ\rq\rq , biến cố B: \lq\lq An lấy được viên bi màu xanh\rq\rq . Hai biến cố A và B là hai
	\choice 
		 {biến cố đối}
		 {biến cố độc lập}
		{\True biến  cố không độc lập}
		{biến cố  hợp}
	\loigiai{ 
	Hai biến cố A và B là hai biến cố không độc lập vì xác suất lấy dược viên bi xanh của An phụ thuộc vào xác suất bạn Hưng lấy được bi màu nào.} 
\end{ex}

% nguồn: Ok_TracNghiemLop11_Lan1
\begin{ex}%[1D9N1-1]%[ma:9eef9f40]
	Cho hai biến cố $A,B$;  biến cố giao của hai biến cố $A,B$ ký hiệu là
	 \choice
		{$A\setminus B$} 
		{$B\setminus A$} 
		{$A\cup  B$}
		{\True  $A\cap B$}
	\loigiai{ 
		Biến cố giao  của hai biến cố $A,B$ ký hiệu là $A\cap B$.
	 }
\end{ex}
